\bigskip
\noindent Autres pistes de l'onglet Findings pour ce thème, dans des cotes qui n'ont pas de section ici :\par

\begin{encadre}
\begin{exemple*}[Grothendieck, pages 30 à 32 et 35, carte des structures spéciales de $M_{0,5}$\piste{148-m05-special-structures-chart}]\nonverifie
Comme points de base et générateurs d'une présentation du groupoïde de Teichmüller de cinq points sur la sphère, construite à partir de pièces de dimension modulaire $\leq 1$, la cote propose une carte de $542$ structures spéciales en neuf orbites sous $\mathfrak S_5$ sur trois niveaux ($12$ pentagonales, $20$ bitétraédriques et $30$ pyramidales ; $60 + 60 + 120$ ; $60 + 120 + 60$), reliées par $1620$ flèches en quatorze orbites.
\end{exemple*}
\end{encadre}
{\small Les relations sont annoncées et jamais écrites. Pour trois types de flèches, les sources et buts dessinés ne correspondent pas aux facteurs du décompte, que l'édition reproduit sans le vérifier, et \q{1620} surcharge un \q{1580} barré. Mathlib n'a pas le groupoïde de Teichmüller. On se reportera aux notes de Grothendieck (\lot{148}{2}{30}).\par}

\begin{encadre}
\begin{definition*}[Grothendieck, pages 40 et 49 à 60, groupoïde de Teichmüller multimarqué\piste{148-multimarked-teichmuller-groupoid}]
Dans le groupoïde de Teichmüller des surfaces découpées, chaque cercle de découpe ou de bord $a$ porte un nombre quelconque $r(a) \geq 0$ de points marqués, qui forment un torseur sous le groupe tordu $\mathbf Z_a/r(a)$. Les opérations élémentaires sont, outre le découpage et le recollement, l'effacement partiel uniforme, restriction au sous-groupe d'ordre $r'$ divisant $r$, et le surmarquage uniforme, extension le long de $\mathbf Z/r \to \mathbf Z/dr$. Le recollement de deux cercles est défini à isomorphisme unique près dès qu'ils portent le même nombre de points et qu'on se donne un anti-isomorphisme de leurs torseurs.
\end{definition*}
\end{encadre}
{\small La définition du surmarquage uniforme et la phrase qui ramène le recollement à $r(a) = r(a')$ contiennent des mots illisibles ; l'unicité du recollement est justifiée par la lecture, et la restriction au sous-groupe d'ordre $r'$ corrige la page. Mathlib n'a pas le groupoïde de Teichmüller des surfaces découpées. On se reportera aux notes de Grothendieck (\lot{148}{2}{40}).\par}
